\chapter*{The Theory at a Glance}
\phantomsection\addcontentsline{toc}{chapter}{The Theory at a Glance}
\markboth{The Theory at a Glance}{}
\label{comp:glance}

\section*{The theory in one page}
\phantomsection\addcontentsline{toc}{section}{The theory in one page}

At one base point $c_*$, a population of agents is a lattice gauge theory on a graph whose site
variables are probability laws. Each agent carries a belief law and a model law on spaces where a
structure group $G$ acts, and a frame saying in which gauge those laws are written; each edge carries
an independent group-valued link, so products of links around loops carry holonomy. Agents interact
through normalized record kernels. From these local data the theory builds one normalized joint law,
one recognition law for the whole population, and one evidence bound, which is the Gibbs variational
principle applied to that joint. Attention is a posterior over which neighbor a receiver listens to,
and selective communication adds query, key, and value maps between the agents' internal spaces.
Fisher geometry turns free-energy descent into oriented histories with an information clock, once a
metric is chosen, and a fast-slow split separates beliefs from models.

Coarse-graining is a Markov kernel from a block of agents to a parent. Applied to the joint, the
posterior, and the recognition law at once, it preserves the evidence and splits the free energy
exactly into the parent's free energy plus the information discarded inside the block. The parent's
law therefore always exists, with an exact effective action. What coarse-graining does not give is a
closed flow on a finite set of couplings, a thermodynamic limit, or a proof that the parent evolves as
an autonomous process. Formation of meta-agents is proved in two settings. In a family with a
hierarchy of interaction strengths supplied by hand, exact parents track a prepared state as it
relaxes through the hierarchy. On a four-agent ring, a pair forms by feedback between attention and
belief relaxation, and slow learning selects one of the two pairings by amplifying an asymmetry of
the initial law. Formation that the agents' coupled inference produces on its own, beyond one level
and beyond four agents, is open.

\section*{Where the chapters stand}
\phantomsection\addcontentsline{toc}{section}{Where the chapters stand}

\colorlet{VCOLch02}{lightgreen}\colorlet{VCOLch03}{lightgreen}\colorlet{VCOLch04}{lightgreen}
\colorlet{VCOLch05}{lightgreen}\colorlet{VCOLch06}{lightgreen}\colorlet{VCOLch07}{lightorange}
\colorlet{VCOLch08}{lightorange}\colorlet{VCOLch09}{lightgreen}\colorlet{VCOLch10}{lightgreen}
\colorlet{VCOLch11}{lightorange}\colorlet{VCOLch12}{lightgreen}\colorlet{VCOLch13}{lightorange}
\colorlet{VCOLch14}{lightorange}\colorlet{VCOLsg}{lightgreen}\colorlet{VCOLsb}{lightorange}
\begin{figure}[htbp]
\centering
\begin{tikzpicture}[
  every node/.style={font=\sffamily\scriptsize,align=center},
  box/.style={draw=primaryblue,rounded corners=2pt,minimum width=3.1cm,minimum height=0.95cm,inner sep=3pt,line width=0.6pt},
  arr/.style={-{Latex[length=2mm]},draw=darkgray,line width=0.5pt},
  sup/.style={box,draw=mediumgray,dashed},
]
\node[box,fill=VCOLch02] (c2) at (0,0) {2\enspace Geometry at $c_*$\\bundles, links, holonomy};
\node[box,fill=VCOLch03] (c3) at (0,-1.35) {3\enspace Probability semantics\\kernels, levels, chain rule};
\node[box,fill=VCOLch04] (c4) at (0,-2.7) {4\enspace Normalized joint law\\typing prohibition};
\node[box,fill=VCOLch05] (c5) at (0,-4.05) {5\enspace One evidence bound\\local--global identity};
\node[box,fill=VCOLch06] (c6) at (4.3,-2.7) {6\enspace Attention and\\selective communication};
\node[box,fill=VCOLch07] (c7) at (4.3,-4.05) {7\enspace Fiber information\\geometry, histories};
\node[box,fill=VCOLch08] (c8) at (4.3,-5.4) {8\enspace Fast beliefs,\\slow models};
\node[box,fill=VCOLch09] (c9) at (0,-6.75) {9\enspace Coarse channels\\exact coarse free energy};
\node[box,fill=VCOLch10] (c10) at (0,-8.1) {10\enspace Network state\\fixed sectors, holonomy};
\node[box,fill=VCOLch11] (c11) at (4.3,-8.1) {11\enspace Exact network\\coarse-graining};
\node[box,fill=VCOLch12] (c12) at (4.3,-9.45) {12\enspace Multiscale VFE,\\closing the tower};
\node[box,fill=VCOLch13] (c13) at (8.6,-9.45) {13\enspace Formation under a\\supplied hierarchy};
\node[box,fill=VCOLch14] (c14) at (8.6,-8.1) {14\enspace Four-cycle selection\\and parent activation};
\node[sup,fill=VCOLsg] (sg) at (8.6,-2.7) {Supplement G\\Gaussian realization};
\node[sup,fill=VCOLsb] (sb) at (8.6,-5.4) {Supplement B\\base, infinite networks,\\lattice cluster};
\draw[arr] (c2) -- (c3);
\draw[arr] (c3) -- (c4);
\draw[arr] (c4) -- (c5);
\draw[arr] (c5) -- (c6);
\draw[arr] (c5) -- (c7);
\draw[arr] (c7) -- (c8);
\draw[arr] (c5) -- (c9);
\draw[arr] (c9) -- (c10);
\draw[arr] (c10) -- (c11);
\draw[arr] (c9) -- (c11);
\draw[arr] (c11) -- (c12);
\draw[arr] (c12) -- (c13);
\draw[arr] (c8) -- (c13);
\draw[arr] (c13) -- (c14);
\draw[arr,dashed,draw=mediumgray] (sg) -- (c6);
\draw[arr,dashed,draw=mediumgray] (sg.south) -- (sb.north);
\draw[arr,dashed,draw=mediumgray] (sb) -- (c11);
\end{tikzpicture}
\caption{How the chapters depend on one another. Fill color gives the companion's verdict on each
technical chapter's own claims: green, delivered as worded; orange, delivered but narrower than the
wording. Dashed boxes are the supplements, on which no main-manuscript result depends. Chapter~1 is
the program statement, whose verdict is ``partly delivered'' (\cref{comp:ch01}), and chapter~15 is
related work.}
\label{comp:fig-map}
\end{figure}

\Cref{comp:fig-map} shows the chapters in dependency order, colored by verdict. No chapter is ``not
delivered'': every chapter proves what its statements say, and the open problems are posed as open.
Where a verdict is ``narrower than it sounds'', the gap is between an opening paragraph or a title
and the theorems, and the companion chapter names the sentence. The theory's unproved ambitions sit
in its open problems, which \hyperref[comp:open-ends]{Open Ends of the Theory} collects; the two
``not delivered'' rows of the table below are two of those ambitions, which the technical documents
themselves pose as open.

\clearpage
\section*{The headline claims}
\phantomsection\addcontentsline{toc}{section}{The headline claims}

\begingroup
\small\renewcommand{\arraystretch}{1.25}
\begin{longtable}{@{}>{\raggedright\arraybackslash}p{0.34\textwidth}>{\raggedright\arraybackslash}p{0.19\textwidth}>{\raggedright\arraybackslash}p{0.39\textwidth}@{}}
\toprule
\textbf{Claim (where)} & \textbf{Verdict} & \textbf{In one sentence}\\
\midrule\endhead
\bottomrule\endfoot
One evidence bound for the whole population, tight exactly at the posterior (\cref{comp:ch05}) & \DELIVERED & The Gibbs variational principle for a correctly normalized joint, written in the extended reals; standard in substance, new in its typing. \\
Local objectives are exact conditional pieces of the collective one (\cref{comp:ch05}) & \DELIVERED & The chain rule for relative entropy with the rest of the population held fixed; the ``exact potential game'' needs singleton blocks with product profiles. \\
Everything is gauge covariant at one base point (\cref{comp:ch02,comp:ch04}) & \DELIVERED & Evidence and bound are invariant under passive frame changes when kernels are covariant; links carry per-edge twist classes and loop holonomy. \\
Attention is inference (\cref{comp:ch06}) & \DELIVERED & The softmax is Bayes' rule for a declared latent pointer; the transformer's free value and output maps are replaced by a tied value map in the exact update. \\
Information geometry and inference histories (\cref{comp:ch07}) & \NARROWERstack & Fisher contraction and the information clock are proved; they cover a meta-agent only when its law is a parameter-independent Markov image, and every orbit rests on a chosen metric. \\
Collective free energy does not factor through marginals (\cref{comp:ch07}) & \DELIVERED & An exact refutation by a three-agent parity witness. \\
Fast beliefs, slow models (\cref{comp:ch08}) & \NARROWERstack & The reduced slow flow is the kinematics of adiabatic elimination; that the actual flow tracks it is proved only on the four-cycle's held sheets. \\
Coarse channels preserve evidence and split the free energy exactly; a parent law always exists (\cref{comp:ch09}) & \DELIVERED & The relative-entropy chain rule for one channel pushing three laws; exact, static, and with no closure or selected channel. \\
Holonomy decides what a coarse description can retain (\cref{comp:ch10}) & \DELIVERED & Fixed sectors and a holonomy-conditioned parent score; for generic $\GL(K)$ links the Gaussian score is infinite on any block with a loop. \\
Exact coarse-graining of the agent network (\cref{comp:ch11}) & \NARROWERstack & The effective action exists for every channel, but closure keeps the complete laws, the coupling calculus excludes Gaussian $\GL(K)$ couplings, everything is finite-network, and there is no thermodynamic limit. \\
One free energy for the multiscale tower; parents influence children (\cref{comp:ch12}) & \DELIVERED & As accounting and typing; parent influence is proved as a typed possibility, not as dynamics. \\
Successive meta-agents form under free-energy descent with learning (\cref{comp:ch13}) & \NARROWERstack & Exact parents track linear relaxation through a hierarchy supplied by hand; attention and learning play no role in the coherence conclusion. \\
A pair forms endogenously on a four-cycle (\cref{comp:ch13}) & \DELIVERED{} \TOY & One pairing episode on four agents, seeded by the random initial assignment, with no learning. \\
Slow learning selects a pairing (\cref{comp:ch14}) & \DELIVERED{} \TOY & A law-level theorem on the four-cycle that amplifies an asymmetry of the initial law; not spontaneous symmetry breaking. \\
Parents activated after formation behave well (\cref{comp:ch14}) & \DELIVERED{} \TOY & On the four-cycle, activation at a deterministic time after the fast modes thermalize works; activation at the first completion is refuted. \\
A parent evolves as an autonomous process (\cref{comp:ch11}) & \NOTDELIVERED & Posed as the defining open problem of the phase, \MT{open:rg-pointwise-parent-dynamics}. \\
Collectives form through the agents' coupled inference, recursively (\cref{comp:ch01,comp:ch13}) & \NOTDELIVERED & The program aim; posed as \MT{open:rg-formation-general} and \MT{open:full-graph-hierarchy}. \\
Gaussian realization: aggregation, Kron reduction, holonomy (\cref{comp:sg1,comp:sg2}) & \DELIVERED & Aggregation is the exact algebraic-multigrid coarse operator; Kron closure is classified; the interaction family is declared, not derived, and complex Gaussians are not treated. \\
Asymptotic renormalization and the lattice cluster (\cref{comp:sb}) & \NARROWERstack & Finite depth gives no asymptotics; the $\mathbb Z_2$ lattice cluster is a genuine constructive result whose sites are not yet typed as agents. \\
\end{longtable}
\endgroup

\section*{What a physicist should take away}
\phantomsection\addcontentsline{toc}{section}{What a physicist should take away}

\paragraph{The foundations are right, and mostly standard in substance.}
The evidence bound, the local objectives, gauge covariance, and exact coarse-graining are proved at
full generality, and most of them are known identities (the Gibbs variational principle, the chain
rule for relative entropy, change of variables, pushforward densities) written in a typed language.
The typing is where the value lies: it separates objects that coincide only in Gaussian or flat
cases, and it prevents collective objects from being sums that merely look global. The one
foundational gap with consequences is that belief-alignment divergences, the attention energies the
program relies on, are law-point energies without a normalized prior over law points
(\MT{open:prob-law-point-prior}), so they are not yet terms of one evidence bound at that level
(\cref{comp:ch03}).

\paragraph{The renormalization at one base point is exact but not closed.}
Every coarse channel produces an exact parent law and an exact effective action. There is no closed
truncation, no flow on a finite set of couplings (a finite-network obstruction rules one out across
reducing steps), no thermodynamic limit, and no autonomous parent dynamics. None of the main
manuscript's open problems takes the limit of many agents, so the thermodynamic-limit view of
renormalization at $c_*$ is not yet posed there (\cref{comp:ch11}, \hyperref[comp:open-ends]{Open
Ends of the Theory}).

\paragraph{Formation is proved only where it is supplied or small.}
The finite-depth theorem shows that exact parents survive a relaxation through a supplied hierarchy
of stiffnesses; it is a statement about tracking, not about emergence. The endogenous results live
on four agents, where selection by learning amplifies an initial asymmetry, as Curie's principle
requires of a unique law-level dynamics. By the small-network rule these results ground the
mechanisms; they neither establish nor exclude formation at scale, and they say nothing for or
against phase transitions, which need the many-agent limit (\cref{comp:ch13,comp:ch14}).

\paragraph{Noncompact groups meet specific friction.}
For the program's target of hundreds of agents with $\GL(K,\R)$, $\GL(K,\mathbb C)$, or $\mathrm{SU}(N)$
links, the documents record several boundaries. A normalized law on noncompact links is not
constructed (\cref{comp:ch02}). The holonomy-conditioned Gaussian parent score is infinite for
generic $\GL(K)$ loops (\cref{comp:ch10}). The coupling calculus of exact network coarse-graining
is proved only in a bounded tier that excludes Gaussian $\GL(K)$ couplings (\cref{comp:ch11}). Gauge
absorption of Gaussian gains is proved on forests and pairs, with only sign conditions on cycles, and
complex Gaussians are not treated (\cref{comp:sg1}). No symmetric rule that ignores a reference point
averages frames in $\GL^+(K)$, so a meta-agent's frame must be chosen with a reference
(\cref{comp:ch10}). $\mathrm{SU}(N)$, being compact, avoids the first two.

\paragraph{What is new.}
Beyond the typing, the new mathematics is concentrated in a few places: the per-edge twist invariants
that a shared fiber gives links at one point (\cref{comp:ch02}), the holonomy-conditioned fixed
sectors and parent score (\cref{comp:ch10}), the parity no-go for factoring collective free energy
(\cref{comp:ch07}), the Kron closure classification (\cref{comp:sg2}), the formation machinery that
carries exact parents through a relaxation (\cref{comp:ch13}), and the four-cycle selection and
activation-clock results (\cref{comp:ch14}).

\section*{How far this companion was checked}
\phantomsection\addcontentsline{toc}{section}{How far this companion was checked}

The companion's chapters were drafted by Claude Opus 5.5 against the technical statements. A second
model, Claude Fable 5.1, checked the reading guide, chapters~1, 5, 6, 9, 10, and~18, and the open
ends against the technical documents, and its findings were applied: none of them changed a verdict,
and several corrected paraphrases that were stronger than their sources. The remaining chapters
(2~to~4, 7, 8, 11~to~17), this overview, and the glossary were not checked by a second model. Their
drafters verified every cited label and recomputed many numerical witnesses, and the central claim
of \cref{comp:ch13}, that attention and learning play no role in the coherence conclusion of the
formation theorem, was checked against its source. Read the unchecked chapters with that in mind,
and trust the technical statement over the companion wherever they differ.
